\documentclass[11pt,a4paper]{article}

\usepackage[utf8]{inputenc}
\usepackage[T2A]{fontenc}
\usepackage[ukrainian,english]{babel}
\usepackage[top=22mm,bottom=22mm,left=22mm,right=22mm]{geometry}
\usepackage{hyperref}
\usepackage{fancyhdr}
\usepackage{parskip}
\usepackage{tcolorbox}
\usepackage{enumitem}
\usepackage{listings}
\usepackage{xcolor}
\usepackage{amssymb}

\tcbuselibrary{skins,breakable}

\hypersetup{
    pdftitle={Workshop 06 --- Plugins: Handout},
    pdfauthor={Vadym Bondarenko},
    colorlinks=true,
    linkcolor=blue,
    urlcolor=cyan,
}

\pagestyle{fancy}
\fancyhf{}
\fancyhead[L]{\small\textit{Workshop 06 --- Plugins}}
\fancyhead[R]{\small\thepage}
\fancyfoot[C]{\small\textit{vcdev.me \textperiodcentered{} 2026}}
\renewcommand{\headrulewidth}{0.4pt}

\lstdefinelanguage{yaml}{
    keywords={true,false,null},
    sensitive=false,
    comment=[l]{\#},
    morestring=[b]',
    morestring=[b]"
}

\lstset{
    basicstyle=\ttfamily\small,
    breaklines=true,
    frame=single,
    backgroundcolor=\color{gray!10},
    rulecolor=\color{gray!40},
    xleftmargin=8pt,
    xrightmargin=8pt,
    aboveskip=6pt,
    belowskip=6pt,
}

\title{%
  {\Huge\bfseries Workshop 06: Plugins}\\[0.6em]
  {\large Глибоке занурення в авторство --- handout}
}
\author{Vadym Bondarenko \\ \small\href{mailto:vadym.bondarenko@vcdev.me}{vadym.bondarenko@vcdev.me}}
\date{2026}

\begin{document}

\maketitle

\begin{abstract}
Пост-воркшоп референс. Не повторює слайди --- сухе зведення схем, маніфестів, чек-листів і pitfalls. Усі факти з офіційних docs (\href{https://code.claude.com/docs/en/plugins}{code.claude.com/docs/en/plugins}, \href{https://code.claude.com/docs/en/plugin-marketplaces}{plugin-marketplaces}, \href{https://code.claude.com/docs/en/plugins-reference}{plugins-reference}).
\end{abstract}

\tableofcontents
\newpage

\section{Що таке плагін}

Самодостатня тека з компонентами, що розширює Claude Code. Маніфест \texttt{.claude-plugin/plugin.json} \textbf{опційний}; без нього Claude Code бере ім'я з теки і авто-знаходить компоненти у дефолтних локаціях.

Компоненти, які можна bundle-ити: \texttt{skills/}, \texttt{commands/}, \texttt{agents/}, \texttt{hooks/}, \texttt{.mcp.json}, \texttt{.lsp.json}, \texttt{monitors/}, \texttt{themes/}, \texttt{output-styles/}, \texttt{bin/}, \texttt{settings.json}.

\textbf{Інвокація:} skill-и і команди --- завжди namespaced: \texttt{/<plugin-name>:<component-name>}. Agents --- у \texttt{/agents} як \texttt{<plugin>:<agent>}.

\section{Plugin manifest cheatsheet}

\subsection{Локація}

Тільки одне правильне місце: \texttt{<plugin-root>/.claude-plugin/plugin.json}.

\textbf{Pitfall:} компоненти (\texttt{skills/}, \texttt{commands/}, \texttt{agents/}, \texttt{hooks/}) --- \textbf{на root плагіна}, НЕ всередині \texttt{.claude-plugin/}. Тільки \texttt{plugin.json} живе у \texttt{.claude-plugin/}.

\subsection{Required і metadata}

\begin{lstlisting}
{
  "name": "git-toolkit",            // REQUIRED. kebab-case
  "version": "0.1.0",               // semver OR omit (commit-SHA)
  "description": "...",
  "author": { "name": "...", "email": "...", "url": "..." },
  "homepage": "https://...",
  "repository": "https://...",
  "license": "MIT",
  "keywords": ["git", "workflow"]
}
\end{lstlisting}

\subsection{Component-path поля}

\begin{lstlisting}
{
  "skills": "./custom/skills/",         // replaces default skills/
  "skills": ["./skills/", "./extras/"], // or array (default + extra)
  "commands": "./custom/commands/",
  "agents": "./custom/agents/",
  "hooks": "./hooks/hooks.json",
  "mcpServers": "./.mcp.json",
  "lspServers": "./.lsp.json",
  "monitors": "./monitors.json",
  "outputStyles": "./styles/",
  "themes": "./themes/"
}
\end{lstlisting}

\textbf{Правила:} усі шляхи відносні до plugin root, починаються з \texttt{./}. Кастомний шлях \textbf{замінює} дефолт. Щоб додати --- масив із \texttt{./skills/} плюс твоє.

\subsection{Просунуті поля}

\begin{lstlisting}
{
  "userConfig": {
    "api_token": {
      "type": "string", "title": "API token",
      "description": "...", "sensitive": true
    }
  },
  "dependencies": [
    { "name": "secrets-vault", "version": "~2.1.0" }
  ]
}
\end{lstlisting}

\texttt{userConfig} --- Claude Code запитає юзера при enable; доступне як \texttt{\$\{user\_config.KEY\}}.

\section{Структура теки плагіна}

\begin{lstlisting}
git-toolkit/
|-- .claude-plugin/
|   `-- plugin.json
|-- skills/
|   `-- status-summary/SKILL.md
|-- commands/
|   `-- log-stats.md
|-- agents/
|   `-- commit-message-reviewer.md
|-- hooks/
|   `-- hooks.json
|-- scripts/
|   `-- log-edit.sh
|-- .mcp.json                (optional)
|-- .lsp.json                (optional)
|-- monitors/monitors.json   (optional)
|-- bin/                     (optional, added to $PATH)
|-- settings.json            (optional, only `agent`, `subagentStatusLine`)
|-- LICENSE
`-- README.md
\end{lstlisting}

\section{Multi-component bundling}

\subsection{Skill (директорія + SKILL.md)}

\begin{lstlisting}
---
description: Use when user asks for current git working tree state
allowed-tools: [Bash]
---
# git-toolkit:status-summary
Run `git status --short --branch` and group by section.
\end{lstlisting}

\subsection{Command (плоский .md)}

\begin{lstlisting}
---
description: Show commit count by author for the last 30 days
allowed-tools: [Bash]
---
Run `git shortlog -sn --since="30 days ago"` and format as table.
\end{lstlisting}

Обидва формати дають \texttt{/git-toolkit:<name>}. Для нових плагінів використовуй \texttt{skills/}; \texttt{commands/} --- legacy.

\subsection{Agent}

\begin{lstlisting}
---
name: commit-message-reviewer
description: Reviews staged commit message
tools: Bash, Read, Grep
model: sonnet
effort: low
---
You are a commit-message reviewer...
\end{lstlisting}

\textbf{Підтримувані поля:} \texttt{name}, \texttt{description}, \texttt{model}, \texttt{effort}, \texttt{maxTurns}, \texttt{tools}, \texttt{disallowedTools}, \texttt{skills}, \texttt{memory}, \texttt{background}, \texttt{isolation} (\texttt{"worktree"} --- єдине валідне).

\textbf{Заборонено для plugin agents:} \texttt{hooks}, \texttt{mcpServers}, \texttt{permissionMode} (security).

\subsection{Hook}

\texttt{hooks/hooks.json}:

\begin{lstlisting}
{
  "hooks": {
    "PostToolUse": [
      {
        "matcher": "Write|Edit",
        "hooks": [
          {
            "type": "command",
            "command": "${CLAUDE_PLUGIN_ROOT}/scripts/log-edit.sh"
          }
        ]
      }
    ]
  }
}
\end{lstlisting}

\textbf{Hook types:} \texttt{command}, \texttt{http}, \texttt{mcp\_tool}, \texttt{prompt}, \texttt{agent}.

\textbf{Найчастіші events:} \texttt{SessionStart}, \texttt{UserPromptSubmit}, \texttt{PreToolUse}, \texttt{PostToolUse}, \texttt{PostToolUseFailure}, \texttt{Stop}, \texttt{SubagentStart}, \texttt{SubagentStop}, \texttt{PreCompact}, \texttt{PostCompact}, \texttt{FileChanged}, \texttt{CwdChanged}, \texttt{InstructionsLoaded}, \texttt{SessionEnd}.

\subsection{Hook input}

Hook отримує JSON \textbf{на stdin} (не як CLI args):

\begin{lstlisting}
#!/usr/bin/env bash
INPUT=$(cat)
FILE=$(echo "$INPUT" | jq -r '.tool_input.file_path')
echo "edited: $FILE" >> "${CLAUDE_PLUGIN_DATA}/edits.log"
\end{lstlisting}

\textbf{Не забудь \texttt{chmod +x}} і shebang.

\section{Plugin path variables}

\begin{tabular}{ll}
\hline
\texttt{\$\{CLAUDE\_PLUGIN\_ROOT\}} & абсолютний шлях до cache-копії плагіна \\
& \emph{змінюється на оновленні; read-only зі сторони plugin-а} \\
\hline
\texttt{\$\{CLAUDE\_PLUGIN\_DATA\}} & персистентна тека стану \\
& \texttt{\textasciitilde/.claude/plugins/data/<id>/} \\
& \emph{виживає update; для логів, кешів, node\_modules} \\
\hline
\end{tabular}

Обидва доступні: inline у skill/agent content, у hook commands, у monitor commands, у MCP/LSP configs. Експортуються як env vars у subprocesses.

\section{Versioning}

\subsection{Resolution order}

Перше, що знайшлось:

\begin{enumerate}
\item \texttt{version} у \texttt{plugin.json}
\item \texttt{version} у marketplace entry
\item Git commit SHA (для \texttt{github}, \texttt{url}, \texttt{git-subdir}, relative path у git-hosted marketplace)
\item \texttt{unknown} (для npm sources або local-dirs не у git)
\end{enumerate}

\subsection{Дві стратегії}

\begin{tabular}{lll}
\hline
\textbf{Стратегія} & \textbf{Як} & \textbf{Update behavior} \\
\hline
Explicit semver & \texttt{"version": "2.1.0"} & Юзери оновлюються лише при bump-і \\
Commit-SHA & Опустити \texttt{version} & Кожен новий commit = нова версія \\
\hline
\end{tabular}

\textbf{Pitfall:} \texttt{version} у двох місцях (\texttt{plugin.json} і marketplace entry) --- \texttt{plugin.json} silently виграє. Лиши одне.

\textbf{Release channels:} два marketplaces на той самий repo з різними \texttt{ref}. Кожен мусить резолвитись у різну версію --- інакше Claude Code вважає однаковими і пропускає update.

\section{Marketplace}

\subsection{Структура}

\begin{lstlisting}
my-marketplace/
|-- .claude-plugin/
|   `-- marketplace.json
`-- plugins/
    `-- git-toolkit/
        |-- .claude-plugin/plugin.json
        |-- skills/, commands/, agents/, hooks/, scripts/
        `-- ...
\end{lstlisting}

\subsection{\texttt{marketplace.json} schema}

\textbf{Required:}

\begin{lstlisting}
{
  "name": "vcdev-marketplace",     // kebab-case, public-facing
  "owner": { "name": "..." },      // email optional
  "plugins": [
    {
      "name": "git-toolkit",       // kebab-case
      "source": "./plugins/git-toolkit"  // or object (see below)
    }
  ]
}
\end{lstlisting}

\textbf{Optional metadata:}

\begin{lstlisting}
{
  "metadata": {
    "description": "...",
    "version": "1.0.0",
    "pluginRoot": "./plugins"   // prefix for relative sources
  },
  "allowCrossMarketplaceDependenciesOn": ["other-marketplace"]
}
\end{lstlisting}

\textbf{Reserved \texttt{name}-и (заборонено):}

\texttt{claude-code-marketplace}, \texttt{claude-code-plugins}, \texttt{claude-plugins-official}, \texttt{anthropic-marketplace}, \texttt{anthropic-plugins}, \texttt{agent-skills}, \texttt{knowledge-work-plugins}, \texttt{life-sciences}.

\subsection{Plugin sources}

\begin{tabular}{ll}
\hline
\textbf{Source} & \textbf{Приклад} \\
\hline
Relative & \texttt{"./plugins/git-toolkit"} \\
GitHub & \texttt{\{ "source": "github", "repo": "owner/repo", "ref"?, "sha"? \}} \\
Git URL & \texttt{\{ "source": "url", "url": "...", "ref"?, "sha"? \}} \\
Git subdir & \texttt{\{ "source": "git-subdir", "url": "...", "path": "...", "ref"? \}} \\
npm & \texttt{\{ "source": "npm", "package": "@org/x", "version"? \}} \\
\hline
\end{tabular}

\textbf{Pitfall:} relative-path source + marketplace додано через raw URL до \texttt{marketplace.json} --- не спрацює (тільки JSON downloaded, repo --- ні). Use \texttt{github}/\texttt{url}/\texttt{git-subdir}/\texttt{npm}.

\section{Installation scopes}

\begin{tabular}{lll}
\hline
\textbf{Scope} & \textbf{Settings file} & \textbf{Use case} \\
\hline
\texttt{user} & \texttt{\textasciitilde/.claude/settings.json} & Особисто, всі проєкти (default) \\
\texttt{project} & \texttt{.claude/settings.json} & Команді через git \\
\texttt{local} & \texttt{.claude/settings.local.json} & Особисто у проєкті, gitignored \\
\texttt{managed} & Managed settings & Org-адміни, read-only \\
\hline
\end{tabular}

\section{CLI команди}

\subsection{Plugin management}

\begin{lstlisting}
claude --plugin-dir ./my-plugin           # local dev (single session)
claude plugin install <name>@<marketplace> [--scope user|project|local]
claude plugin uninstall <name> [--keep-data]
claude plugin enable <name> / plugin disable <name>
claude plugin update <name>
claude plugin list [--json] [--available]
claude plugin validate .                  # check JSON + frontmatter
claude plugin tag [--push] [--dry-run]    # release git tag
\end{lstlisting}

\subsection{Marketplace management}

\begin{lstlisting}
claude plugin marketplace add <source> [--scope ...] [--sparse paths...]
claude plugin marketplace list
claude plugin marketplace remove <name>
claude plugin marketplace update [name]
\end{lstlisting}

В інтерактивній сесії --- ті самі через \texttt{/plugin marketplace add}, \texttt{/plugin install}, \texttt{/plugin validate}, \texttt{/reload-plugins}.

\section{Pitfalls --- 10 граблів}

\begin{enumerate}
\item Компоненти всередині \texttt{.claude-plugin/} --- мають бути на root.
\item Set \texttt{version} і не bump-нути --- юзери не отримують update.
\item Version у двох місцях (\texttt{plugin.json} + marketplace entry) --- \texttt{plugin.json} виграє.
\item Hard-coded paths замість \texttt{\$\{CLAUDE\_PLUGIN\_ROOT\}} --- silent fail після install.
\item Path-traversal \texttt{../shared-utils} --- не копіюється у cache.
\item Hook script не executable --- \texttt{chmod +x} + shebang.
\item Marketplace через raw URL + relative-path source --- репо не клонується.
\item Plugin agent з \texttt{hooks}/\texttt{mcpServers}/\texttt{permissionMode} --- silently rejected.
\item Marketplace \texttt{name} --- reserved (\texttt{claude-plugins-official} etc.).
\item Plugin не з'являється --- спершу \texttt{claude plugin validate .}, потім \texttt{claude --debug}.
\end{enumerate}

\section{Production-чек-ліст}

\begin{itemize}[leftmargin=2em]
\item[$\square$] \texttt{name} kebab-case, не reserved
\item[$\square$] \texttt{description}, \texttt{author}, \texttt{license} заповнені
\item[$\square$] Стратегія \texttt{version} обрана: explicit АБО omit
\item[$\square$] \texttt{.claude-plugin/} містить ТІЛЬКИ \texttt{plugin.json}
\item[$\square$] \texttt{\$\{CLAUDE\_PLUGIN\_ROOT\}} у hook commands і MCP/LSP configs
\item[$\square$] \texttt{\$\{CLAUDE\_PLUGIN\_DATA\}} для логів, кешів, state
\item[$\square$] \texttt{chmod +x} + shebang на hook scripts
\item[$\square$] Жодного \texttt{..} у paths
\item[$\square$] \texttt{README.md} --- install instructions
\item[$\square$] \texttt{claude plugin validate .} проходить
\item[$\square$] Тестовано через marketplace flow, не лише \texttt{--plugin-dir}
\end{itemize}

\section{Команди для вправ}

\subsection{Вправа 1 --- ініціалізація}

\begin{lstlisting}
mkdir -p git-toolkit/.claude-plugin
mkdir -p git-toolkit/{skills,commands,agents,hooks}
# create plugin.json with name, version, description, author
claude plugin validate ./git-toolkit
claude --plugin-dir ./git-toolkit
\end{lstlisting}

\subsection{Вправа 2 --- multi-component}

\begin{lstlisting}
# skill (directory with SKILL.md)
mkdir -p git-toolkit/skills/status-summary
# command (flat .md)
# agent (flat .md with tools/model)
claude plugin validate ./git-toolkit
\end{lstlisting}

\subsection{Вправа 3 --- hook}

\begin{lstlisting}
# hooks/hooks.json: PostToolUse(Write|Edit) -> ${CLAUDE_PLUGIN_ROOT}/scripts/log-edit.sh
# script appends to ${CLAUDE_PLUGIN_DATA}/edits.log
chmod +x git-toolkit/scripts/log-edit.sh
\end{lstlisting}

\subsection{Вправа 4 --- marketplace}

\begin{lstlisting}
# vcdev-marketplace/.claude-plugin/marketplace.json
mv git-toolkit vcdev-marketplace/plugins/
claude plugin validate ./vcdev-marketplace
# in session:
/plugin marketplace add ./vcdev-marketplace
/plugin install git-toolkit@vcdev-marketplace
/git-toolkit:status-summary
\end{lstlisting}

\section{Resources}

\begin{itemize}
\item \href{https://code.claude.com/docs/en/plugins}{code.claude.com/docs/en/plugins} --- Create plugins
\item \href{https://code.claude.com/docs/en/plugin-marketplaces}{code.claude.com/docs/en/plugin-marketplaces} --- Marketplaces
\item \href{https://code.claude.com/docs/en/plugins-reference}{code.claude.com/docs/en/plugins-reference} --- Complete reference
\item \href{https://code.claude.com/docs/en/skills}{code.claude.com/docs/en/skills} --- Skills (workshop 04)
\item Цей репо: \texttt{workshops/06-plugins/exercises/} і \texttt{solutions/}
\end{itemize}

\end{document}
